\mxunit{premiere-c03}{Calcul vectoriel et produit scalaire}\label{C03}
\begin{keypoint}[Le parcours du chapitre]
\textbf{Objectifs.} Choisir entre projection, coordonnées et angle ; calculer une longueur ; caractériser une orthogonalité et un lieu.

\textbf{Prérequis.} Vecteurs, Chasles, Pythagore, cosinus dans un triangle rectangle, repère orthonormé.

\textbf{Repère annuel.} 6–15 octobre : 8 h\end{keypoint}
\section*{Diagnostic : \mxcount{diag} questions, 10 minutes}\label{C03-D}
Répondre sans consulter le cours. Indiquer les questions pour lesquelles une aide est nécessaire. Ce diagnostic ne donne pas lieu à une note.
\begin{enumerate}
\item Pour $A\pt{1;2}$ et $B\pt{4;-2}$, calculer $\vect{AB}$.
\item Calculer la norme de $\pt{3;-4}$.
\item Simplifier $\vect{AB}+\vect{BC}$.
\item Développer $(a+b)^2$.
\item Dans un triangle rectangle, un angle a pour côtés adjacent 3 et hypoténuse 5. Donner son cosinus.
\item Le milieu de $A\pt{-2;1}$ et $B\pt{4;3}$ ?
\mxcounted{diag}
\end{enumerate}
\paragraph{Après le diagnostic.} Reprendre les notions signalées avec le professeur ; refaire une question analogue, puis expliquer la règle utilisée.
\paragraph{Oral et logique.} Comparer, sur un même problème, coordonnées et calcul vectoriel sans repère. Justifier pourquoi votre méthode utilise les données les plus directement.
\mxnewpage
\section*{Chercher avant de formaliser}\label{C03-RECH}
Construire trois points M tels que $\vect{MA}\cdot\vect{MB}=0$ pour un segment AB de 6 cm. Conjecturer le lieu puis produire une preuve qui inclut A et B.
\paragraph{Organisation.} Recherche individuelle 5 min, confrontation à deux 5 min, mise en commun 5 min. Garder une trace d’un essai, même non abouti. Une conjecture doit ensuite être justifiée.
\mxnewpage
\section{Projection et angle}\label{C03-U1}
\begin{keypoint}[Définition et résultats]
Pour deux vecteurs non nuls $\vec u,\vec v$, leur produit scalaire est le réel $\vec u\cdot\vec v=\norm{\vec u}\norm{\vec v}\cos\theta$, où $\theta\in\intff{0}{\pi}$ est leur angle non orienté. Avec un vecteur nul, le produit vaut 0. Si $H$ est le projeté de B sur $(OA)$, $\vect{OA}\cdot\vect{OB}$ vaut $OA\times OH$ lorsque H est sur la demi-droite $[OA)$, et son opposé sinon.
\end{keypoint}
\paragraph{Statut et justification.} Définition ; l’équivalence avec la projection découle du cosinus dans les triangles rectangles, avec une projection algébrique négative pour un angle obtus.
\begin{example}Pour deux vecteurs de normes 3 et 4 formant $60^\circ$, le produit vaut $\dfrac{3\times 4\times 1}{2}=6$. Avec $120^\circ$, il vaut $-6$. Si $A\pt{3;0}$ et $B\pt{-2;4}$, le projeté est $H\pt{-2;0}$ et le produit vaut $-6$.\end{example}\label{C03-EX1}
\mxinclude{figures/C03-U1}
\begin{pitfall}Le résultat est un nombre, pas un vecteur. L’angle droit correspond à 0, pas à 1.\end{pitfall}
\begin{application}Deux vecteurs ont pour normes 2 et 5 et forment $90^\circ$, puis $180^\circ$. Calculer les produits.\end{application}\label{C03-AP1}
\section{Coordonnées, norme et orthogonalité}\label{C03-U2}
\begin{keypoint}[Définition et résultats]
Dans une base orthonormée, pour $\vec u=\pt{x;y}$ et $\vec v=\pt{x\prime;y\prime}$, $\vec u\cdot\vec v=xx\prime+yy\prime$ et $\norm{\vec u}^2=x^2+y^2$. Les vecteurs sont orthogonaux si et seulement si leur produit est nul. Pour une base $\pt{\vec i;\vec j}$ orthonormée, $x=\vec u\cdot\vec i$ et $y=\vec u\cdot\vec j$.
\end{keypoint}
\paragraph{Statut et justification.} Propriétés démontrées à l’aide de la bilinéarité : développer $(x\vec i+y\vec j)\cdot(x\prime\vec i+y\prime\vec j)$ et utiliser $\vec i\cdot\vec i=\vec j\cdot\vec j=1$, $\vec i\cdot\vec j=0$.
\begin{example}$\vec u=\pt{2;3}$, $\vec v=\pt{3;-2}$ : produit $6-6=0$. Les deux vecteurs sont orthogonaux. Leur norme vaut $\sqrt{13}$. Le produit de $\pt{2;3}$ avec $\pt{1;0}$ vaut 2 : on retrouve l’abscisse.\end{example}\label{C03-EX2}
\begin{pitfall}La formule usuelle suppose une base orthonormée. Colinéarité et orthogonalité sont des propriétés différentes.\end{pitfall}
\begin{application}Pour quel réel t les vecteurs $\pt{t;2}$ et $\pt{3;-6}$ sont-ils orthogonaux ?\end{application}\label{C03-AP2}
\section{Calcul vectoriel et identités}\label{C03-U3}
\begin{keypoint}[Définition et résultats]
Le produit scalaire est symétrique et bilinéaire : $\vec u\cdot(\vec v+\vec w)=\vec u\cdot\vec v+\vec u\cdot\vec w$ et $(\lambda\vec u)\cdot\vec v=\lambda(\vec u\cdot\vec v)$. Ainsi $\norm{\vec u\pm\vec v}^2=\norm{\vec u}^2\pm2\vec u\cdot\vec v+\norm{\vec v}^2$. On en tire $\vec u\cdot \vec v=\dfrac{\norm{\vec u+\vec v}^2-\norm{\vec u}^2-\norm{\vec v}^2}{2}$.
\end{keypoint}
\paragraph{Statut et justification.} Bilinéarité admise à partir de la définition géométrique ; symétrie issue du cosinus. Les identités sont démontrées en développant $(\vec u\pm\vec v)\cdot(\vec u\pm\vec v)$.
\begin{example}Si $\norm{\vec u}=3$, $\norm{\vec v}=4$ et $\norm{\vec u+\vec v}=5$, le produit vaut $\dfrac{25-9-16}{2}=0$. On déduit l’orthogonalité sans coordonnées.\end{example}\label{C03-EX3}
\begin{pitfall}En général, $\norm{\vec u+\vec v}\ne\norm{\vec u}+\norm{\vec v}$. Les identités concernent les carrés des normes.\end{pitfall}
\begin{application}On donne $\norm{\vec u}=2$, $\norm{\vec v}=3$, $\vec u\cdot\vec v=1$. Calculer $\norm{\vec u-\vec v}$.\end{application}\label{C03-AP3}
\section{Al-Kashi et choix d’une méthode}\label{C03-U4}
\begin{keypoint}[Définition et résultats]
Dans un triangle ABC, $BC^2=AB^2+AC^2-2AB\,AC\cos\widehat{BAC}$. Cette formule permet de calculer le troisième côté à partir de deux côtés et de l’angle compris, ou de déterminer le cosinus d’un angle à partir des trois côtés.
\end{keypoint}
\paragraph{Statut et justification.} Par Chasles, $\vect{BC}=\vect{AC}-\vect{AB}$. Développer son carré scalaire donne $BC^2=AC^2+AB^2-2\vect{AB}\cdot\vect{AC}$ ; introduire l’expression avec le cosinus.
\begin{example}$AB=5$, $AC=3$ et $\widehat{BAC}=60^\circ$ : $BC^2=25+9-15=19$, donc $BC=\sqrt{19}$. Avec trois côtés, isoler le cosinus avant d’approcher l’angle à la calculatrice si elle est autorisée.\end{example}\label{C03-EX4}
\begin{pitfall}L’angle est compris entre les deux côtés du terme produit. Le signe moins est essentiel.\end{pitfall}
\begin{application}Un triangle a $AB=AC=4$ et $BC=4\sqrt2$. Déterminer l’angle A.\end{application}\label{C03-AP4}
\section{Le cercle de diamètre AB}\label{C03-U5}
\begin{keypoint}[Définition et résultats]
Pour deux points distincts A et B de milieu I, $\vect {MA}\cdot \vect {MB}=MI^2-\dfrac{AB^2}{4}$. L’ensemble des points M tels que $\vect{MA}\cdot\vect{MB}=0$ est le cercle de centre I et de rayon $\dfrac{AB}{2}$, points A et B compris.
\end{keypoint}
\paragraph{Statut et justification.} Écrire $\vect{MA}=\vect{MI}+\vect{IA}$ et $\vect{MB}=\vect{MI}-\vect{IA}$. Les termes croisés s’annulent ; il reste $MI^2-IA^2$. Comme $IA=\dfrac{AB}{2}$, l’égalité à zéro équivaut à l’équation métrique du cercle.
\begin{example}$A\pt{-2;0}$ et $B\pt{2;0}$ : I est l’origine. Le lieu est $x^2+y^2=4$. Pour $M\pt{0;1}$, le produit vaut $1-4=-3$ ; pour $M\pt{0;2}$, il vaut 0.\end{example}\label{C03-EX5}
\mxinclude{figures/C03-U5}
\begin{pitfall}Les extrémités appartiennent au lieu même si « triangle rectangle » ne convient pas à un triangle dégénéré.\end{pitfall}
\begin{application}Pour $A\pt{0;0}$, $B\pt{6;0}$, déterminer les M tels que $\vect{MA}\cdot\vect{MB}=7$.\end{application}\label{C03-AP5}
\mxnewpage
\section*{Parcours d’exercices}\label{C03-PARC}
Commencer par les exercices 1 et 2. Choisir ensuite un exercice 3 ou 4 sans aide. Les exercices 5 à 7 sont une réserve de transfert, de recherche ou de programmation : le professeur en désigne un. L’exercice 8 se reprend à distance.
\begingroup
\mxexercisehead{C03-E1}{Calculs directs}
Dans une base orthonormée du plan, on considère les vecteurs $\vec u=\pt{1;-2}$ et $\vec v=\pt{4;3}$.
\begin{enumerate}[label=\arabic*.]
\item Calculer le produit scalaire $\vec u\cdot\vec v$, puis les normes $\norm{\vec u}$ et $\norm{\vec v}$.
\item Calculer les coordonnées de $\vec u+\vec v$ et en déduire $\norm{\vec u+\vec v}^2$.
\item Retrouver ce dernier résultat à l’aide de l’identité faisant intervenir $\vec u\cdot\vec v$.
\end{enumerate}
\mxexercisehead{C03-E2}{Un triangle rectangle ?}
Le plan est muni d’un repère orthonormé. On donne les points $A\pt{1;1}$, $B\pt{4;2}$ et $C\pt{0;4}$.
\begin{enumerate}[label=\arabic*.]
\item Calculer les coordonnées de $\vect{AB}$ et $\vect{AC}$, puis leur produit scalaire.
\item En déduire que le triangle $ABC$ est rectangle ; préciser en quel sommet.
\item Calculer les longueurs des deux côtés perpendiculaires, puis l’aire du triangle en unités d’aire.
\end{enumerate}
\mxexercisehead{C03-E3}{Longueur sans repère}
Le quadrilatère $ABCD$ est un parallélogramme tel que $AB=4$, $AD=3$ et $\widehat{BAD}=60^\circ$. Les longueurs sont exprimées dans la même unité.
\begin{enumerate}[label=\arabic*.]
\item Faire une figure et calculer le produit scalaire $\vect{AB}\cdot\vect{AD}$.
\item Déterminer les longueurs exactes des diagonales $[AC]$ et $[BD]$. Expliquer comment les vecteurs correspondant aux diagonales s’expriment à partir de $\vect{AB}$ et $\vect{AD}$.
\item Indiquer quelle diagonale est la plus longue et vérifier que ce résultat est cohérent avec votre figure.
\end{enumerate}
\mxexercisehead{C03-E4}{Retrouver un angle}
Deux vecteurs non nuls $\vec u$ et $\vec v$ vérifient $\norm{\vec u}=2$, $\norm{\vec v}=6$ et $\vec u\cdot\vec v=-6$. On note $\theta$ leur angle non orienté, compris entre $0^\circ$ et $180^\circ$.
\begin{enumerate}[label=\arabic*.]
\item Déterminer la valeur exacte de $\cos\theta$, puis la mesure de $\theta$.
\item Expliquer en quoi le signe du produit scalaire permettait de prévoir que l’angle serait obtus.
\end{enumerate}
\mxexercisehead{C03-E5}{Un lieu variable}
Les points $A$ et $B$ sont distincts et fixés. On note $I$ le milieu de $[AB]$. Pour un réel $k$ fixé, on cherche l’ensemble des points $M$ du plan qui vérifient $\vect{MA}\cdot\vect{MB}=k$.
\begin{enumerate}[label=\arabic*.]
\item Exprimer ce produit scalaire en fonction de $MI$ et de $AB$ en utilisant la relation de Chasles et une identité scalaire.
\item Décrire précisément l’ensemble recherché selon la valeur de $k$. Pour chaque cas, indiquer s’il est vide, réduit à un point ou constitué d’un cercle ; donner alors ses caractéristiques.
\item Expliquer le rôle du fait qu’une distance au carré est positive ou nulle.
\end{enumerate}
\mxexercisehead{C03-E6}{Travail d’une force}
Un objet se déplace en ligne droite de $A$ vers $B$ sur une distance de 4 m. Il est soumis à une force constante $\vec F$ de norme 50 N. Le travail de cette force est défini par $W=\vec F\cdot\vect{AB}$ et s’exprime en joules.
\begin{enumerate}[label=\arabic*.]
\item Calculer $W$ lorsque l’angle entre $\vec F$ et $\vect{AB}$ vaut $60^\circ$.
\item Reprendre le calcul pour un angle de $120^\circ$.
\item Comparer les signes des deux travaux et expliquer, dans chaque situation, si la force favorise le déplacement ou s’y oppose.
\end{enumerate}
\mxexercisehead{C03-E7}{Contraposée}
On souhaite démontrer l’affirmation suivante, valable pour deux vecteurs du plan :\[\norm{\vec u+\vec v}^2\ne\norm{\vec u}^2+\norm{\vec v}^2\quad\Longrightarrow\quad \vec u\text{ et }\vec v\text{ ne sont pas orthogonaux}.\]
\begin{enumerate}[label=\arabic*.]
\item Énoncer la contraposée de cette implication.
\item Démontrer cette contraposée à l’aide d’une identité scalaire, puis conclure sur l’affirmation initiale.
\end{enumerate}
\mxexercisehead{C03-E8}{Reprise différée}
Un triangle $ABC$ vérifie $AB=3$, $AC=4$ et $BC=5$. On cherche $\vect{AB}\cdot\vect{AC}$ sans utiliser de coordonnées.
\begin{enumerate}[label=\arabic*.]
\item Calculer ce produit scalaire par une formule faisant intervenir les trois longueurs.
\item Retrouver le résultat en déterminant la nature du triangle. Citer le théorème utilisé et son hypothèse numérique.
\end{enumerate}
\paragraph{Travail personnel.} Effectuer uniquement la tâche désignée à l’issue de la séance, dans la progression datée : 10–20 min maximum. Les exercices dépendent des notions déjà étudiées ; certaines questions sont donc différées. Reprendre l’exercice 8 une semaine après le chapitre (10 min). Si le temps prévu est dépassé, noter la difficulté et apporter les essais ; ne pas accumuler toute la réserve.
\endgroup
\mxneed{12}\section*{À faire dans le manuel CQFD}
Ces exercices se traitent dans votre manuel. La progression indique la séance ; ils remplacent le travail personnel prévu ce jour-là (10 à 20 minutes au total).
\paragraph{CQFD, exercice 9 p. 250, questions a et b.}\label{C03-M1} Calculer le produit scalaire dans la base orthonormée et détailler les deux produits de coordonnées.
\paragraph{CQFD, exercice 11 p. 250, questions 1, 2 et 3.}\label{C03-M2} Utiliser successivement les coordonnées, la relation de Chasles et le fait que I est le milieu.
\begingroup\small
\mxnewpage\section*{QCM 1 : comprendre et choisir}\label{C03-Q1}
\mxCount{qcm1} questions pédagogiques à réponse unique. Choisir une réponse par question ; expliquer ensuite un distracteur. Ce format de chapitre ne reproduit pas le format officiel de l’EAM.
\begin{enumerate}
\item $\pt{2;1}\cdot\pt{3;-4}$ vaut
\begin{enumerate}[label=\Alph*.,itemsep=2pt]
\item $2$
\item $10$
\item $\pt{6;-4}$
\item $-2$
\end{enumerate}
\item La norme de $\pt{-3;4}$ vaut
\begin{enumerate}[label=\Alph*.,itemsep=2pt]
\item $7$
\item $25$
\item $5$
\item $1$
\end{enumerate}
\item Deux vecteurs de produit nul sont
\begin{enumerate}[label=\Alph*.,itemsep=2pt]
\item égaux
\item colinéaires
\item opposés
\item orthogonaux
\end{enumerate}
\item Normes 2 et 3, angle $60^\circ$ : produit
\begin{enumerate}[label=\Alph*.,itemsep=2pt]
\item $6$
\item $3$
\item $\sqrt3$
\item $0$
\end{enumerate}
\item $\norm{u-v}^2$ vaut
\begin{enumerate}[label=\Alph*.,itemsep=2pt]
\item $\norm u^2-\norm v^2$
\item $\norm u^2+\norm v^2$
\item $\norm u^2-2u\cdot v+\norm v^2$
\item $\norm u^2+2u\cdot v+\norm v^2$
\end{enumerate}
\mxcounted{qcm1}
\end{enumerate}
\endgroup
\begingroup\small
\mxneed{18}\section*{QCM 2 : comprendre et choisir}\label{C03-Q2}
\mxCount{qcm2} questions pédagogiques à réponse unique. Choisir une réponse par question ; expliquer ensuite un distracteur. Ce format de chapitre ne reproduit pas le format officiel de l’EAM.
\begin{enumerate}
\item Un angle obtus entre vecteurs non nuls donne un produit
\begin{enumerate}[label=\Alph*.,itemsep=2pt]
\item positif
\item nul
\item égal à 1
\item négatif
\end{enumerate}
\item Dans une base orthonormée, $u\cdot i$ est
\begin{enumerate}[label=\Alph*.,itemsep=2pt]
\item l’abscisse de u
\item la norme de u
\item l’ordonnée de u
\item toujours nul
\end{enumerate}
\item Le lieu $\vect{MA}\cdot\vect{MB}=0$ pour $A\ne B$ est
\begin{enumerate}[label=\Alph*.,itemsep=2pt]
\item la médiatrice
\item la droite AB
\item le cercle de diamètre AB
\item le milieu seul
\end{enumerate}
\item Si $u\cdot v=3$, $(2u)\cdot(-v)$ vaut
\begin{enumerate}[label=\Alph*.,itemsep=2pt]
\item $-3$
\item $-6$
\item $6$
\item $9$
\end{enumerate}
\item Pour $AB=AC=2$ et angle A droit, $BC$ vaut
\begin{enumerate}[label=\Alph*.,itemsep=2pt]
\item $4$
\item $2$
\item $8$
\item $2\sqrt2$
\end{enumerate}
\mxcounted{qcm2}
\end{enumerate}
\subsection*{Vrai ou faux ? Justifier.}\label{C03-VF}
\begin{enumerate}
\item Le produit scalaire de deux vecteurs est toujours positif.
\item Le vecteur nul est orthogonal à tout vecteur.
\item $u\cdot v=0$ implique que l’un des vecteurs est nul.
\end{enumerate}
\begin{keypoint}[Bilan personnel]Je distingue norme et carré de norme ; je choisis une formule adaptée ; je vérifie le repère ; je rédige une équivalence pour un lieu.

Pour chaque objectif, indiquer : acquis sans aide ; réussi avec aide ; à retravailler. Écrire un exemple et une erreur à éviter.\end{keypoint}
\endgroup
\ifmxteacher\mxinclude{corriges/C03}\fi
